\section{Conclusiones}
El desarrollo de esta práctica se llevó a cabo gracias al tutorial de raspberry pi de la facultad \cite{FI_UNAM_Raspberry_2026} y a las clases del profesor Moisés \cite{Melendez_Reyes_2026}.

\subsection{Hernandez Acosta Mauricio Gabriel}
Durante el desarrollo de esta práctica, logré consolidar mis conocimientos sobre el manejo directo de la memoria en la arquitectura ARM. La resolución de los ejercicios me permitió comprobar la importancia de estructurar correctamente los ciclos y las ramificaciones condicionales para procesar arreglos y matrices sin errores lógicos. El uso del direccionamiento indexado con desplazamiento (offset) resultó ser una herramienta clave para recorrer y manipular bloques de datos de 32 bits de manera precisa. En conclusión, esta experiencia fue fundamental para entender que la programación a bajo nivel requiere una orquestación cuidadosa de los registros para que actúen como punteros efectivos.

\subsection{Perez Olvera Alexis Abraham}
La realización de estas actividades evidenció la enorme ventaja que ofrecen los modos de direccionamiento indirecto para la optimización del código en lenguaje ensamblador. Al programar los algoritmos de búsqueda y ordenamiento, comprobé que técnicas como el direccionamiento post-indexado y los desplazamientos lógicos (LSL) reducen drásticamente la cantidad de instrucciones necesarias para iterar sobre grandes estructuras de datos. Entender cómo el procesador gestiona y actualiza las direcciones de memoria de forma automática me brindó la perspectiva necesaria para escribir rutinas que no solo sean funcionales, sino también compactas y eficientes en su ejecución.

\subsection{Sierra Garcia Mariana}
Esta práctica representó un excelente ejercicio para trasladar la teoría de estructuras de datos a su aplicación real a nivel de hardware. Pude asimilar cómo operaciones que parecen complejas, tales como la multiplicación de matrices, la generación de secuencias o el ordenamiento de elementos, se resuelven netamente mediante la aritmética de punteros. Emplear las distintas variantes de direccionamiento y desplazamientos con LSL me ayudó a dominar el acceso secuencial a la memoria. En resumen, comprendí que la correcta administración y actualización de los registros es la base para interactuar exitosamente con cualquier bloque de información en el sistema.